\section{TASK 2: TOKENIZATION} \label{sec:task2}

We compare \textbf{BERT-base uncased}, trained primarily on natural-language
text, and \textbf{UniXcoder-base}, whose pre-training also includes source
code. Both tokenizers are evaluated with a maximum sequence length of 512.

\subsection{Tokenization of SSH Commands} \label{subsec:task2_commands}

\textbf{Command-level comparison.} Table~\ref{tab:task2_commands} shows the
subword decomposition of the nine required commands. Prefix markers used by the
tokenizers are omitted from the table for readability. BERT uses
\texttt{\#\#} for continuation pieces, whereas UniXcoder uses a spacing marker
on word-initial tokens.

BERT produces 16 tokens for the nine words, a token-to-word ratio of
\textbf{1.78}; UniXcoder produces 14, yielding the better ratio of
\textbf{1.56}. Common words such as \texttt{cat}, \texttt{echo}, and
\texttt{shell} are retained by both vocabularies because they are sufficiently
frequent in their pre-training corpora. UniXcoder additionally preserves
\texttt{grep}, consistent with its exposure to code and command-line text.

\subsection{Full-Corpus Tokenization and Context Limit}
\label{subsec:task2_corpus}

\textbf{Token counts before preprocessing.} Table~\ref{tab:task2_raw_stats}
summarizes tokenization of all 251 training sessions. Although UniXcoder has a
lower median, its mean and maximum are much larger because rare shell strings
are decomposed rather than collapsed into unknown tokens. Across the corpus,
BERT emits \textbf{303 \texttt{[UNK]} tokens}, whereas UniXcoder emits none.

\begin{table}[tbp]
    \begin{minipage}[t]{0.49\linewidth}
        \centering
        \small
        \begin{tabular}{lll}
            \toprule
            \textbf{Command} & \textbf{BERT} & \textbf{UniXcoder} \\
            \midrule
            \texttt{cat}      & cat          & cat \\
            \texttt{shell}    & shell        & shell \\
            \texttt{echo}     & echo         & echo \\
            \texttt{top}      & top          & top \\
            \texttt{chpasswd} & ch, pass, wd & ch, passwd \\
            \texttt{crontab}  & cr, ont, ab  & c, ront, ab \\
            \texttt{wget}     & w, get       & w, get \\
            \texttt{busybox}  & busy, box    & busy, box \\
            \texttt{grep}     & gr, ep       & grep \\
            \bottomrule
        \end{tabular}
        \vspace{0.1cm}
        \caption{Subword tokenization of the SSH commands.}
        \label{tab:task2_commands}
    \end{minipage}
    \hfill
    \begin{minipage}[t]{0.49\linewidth}
        \centering
        \small
        \begin{tabular}{lrr}
            \toprule
            \textbf{Metric} & \textbf{BERT} & \textbf{UniXcoder} \\
            \midrule
            Mean tokens/session       & 178.59 & 409.32 \\
            Median tokens/session     & 71     & 62 \\
            Maximum tokens/session    & 1,889  & 28,920 \\
            Corpus \texttt{[UNK]} count & 303    & 0 \\
            Sessions above 512        & 24     & 29 \\
            \bottomrule
        \end{tabular}
        \vspace{0.1cm}
        \caption{Training-corpus tokenization before word truncation.}
        \label{tab:task2_raw_stats}
    \end{minipage}
\end{table}

Consequently, \textbf{24 BERT sessions (9.56\%)} and \textbf{29 UniXcoder
sessions (11.55\%)} exceed 512 tokens. Figure~\ref{fig:task2_token_ecdf} shows
the corresponding ECDFs and the context threshold. Most sessions remain below
the limit, but UniXcoder has a visibly longer tail.

\begin{figure}[tbp]
    \centering
    \includegraphics[width=\linewidth]{\imagespath Task2/token_count_ecdf.png}
    \caption{ECDF of token counts before preprocessing. The dashed line marks
    the 512-token context limit.}
    \label{fig:task2_token_ecdf}
\end{figure}

\vspace{0.2em}

\noindent
\textbf{Longest session.} Both maxima occur on the same \textbf{134-word}
training session. Two words exceed 30 characters; the longest is a
\textbf{46,860-character encoded payload}, which is kept strictly as text.
BERT replaces portions of the rare sequence with \texttt{[UNK]} and produces
1,889 tokens, whereas UniXcoder decomposes it into 28,920 tokens. The
code-oriented tokenizer is therefore compact on ordinary commands but much
more verbose on pathological strings.

\subsection{Truncating Long Bash Words} \label{subsec:task2_word_truncation}

\textbf{Preprocessing effect.} Following the prescribed preprocessing, every
Bash word longer than 30 characters is shortened before tokenization. This
reduces the average sequence length to \textbf{128.38 tokens for BERT} and
\textbf{110.52 for UniXcoder}; the maxima fall to 920 and 824, respectively.
Only \textbf{six sessions per tokenizer} remain above the 512-token limit.

Figure~\ref{fig:task2_words_tokens} relates Bash-word counts to token counts
after preprocessing. Using the notebook's per-session definition
\(\frac{1}{N}\sum_i \frac{\text{tokens}_i}{\text{words}_i}\), the average
token-to-word ratio is \textbf{3.23 for BERT} and \textbf{2.77 for
UniXcoder}. Thus, once abnormally long words are controlled, UniXcoder provides
the more compact representation of typical SSH syntax.

\begin{figure}[tbp]
    \centering
    \includegraphics[width=\linewidth]{\imagespath Task2/words_vs_tokens.png}
    \caption{Number of Bash words versus generated tokens after truncating
    individual words to 30 characters.}
    \label{fig:task2_words_tokens}
\end{figure}
